% Source: https://artofproblemsolving.com/community/c3h1396472_geometric_sequence
% Stable identity: aops-topic-1396472-q01
\textbf{(a)} With common ratio $r>0$, $a_8=a_2r^6$, so $r^6=\dfrac{320}{5}=64$ and $r=2$. Then $a_4=a_2r^2=5\cdot4=\boxed{20}$.

\textbf{(b)} Let $a_j=a_1r^{j-1}$ with $a_1,r>0$. Then $\dfrac{a_n}{a_m}=r^{n-m}$, so $r=\left(\dfrac{a_n}{a_m}\right)^{\dfrac{1}{n-m}}$ and
\[
a_k=a_m r^{k-m}=a_m\left(\dfrac{a_n}{a_m}\right)^{\dfrac{k-m}{n-m}}=a_m^{1-\dfrac{k-m}{n-m}}\,a_n^{\dfrac{k-m}{n-m}}=a_m^{\dfrac{n-k}{n-m}}\,a_n^{\dfrac{k-m}{n-m}} .
\]
The exponents are the fractions of the way from $n$ back to $k$ and from $m$ to $k$; they add up to $1$. In (a): $a_4=5^{\dfrac{4}{6}}\cdot320^{\dfrac{2}{6}}=\left(5^2\cdot320\right)^{\dfrac{1}{3}}=8000^{\dfrac{1}{3}}=20$. (Taking logarithms, $\log a_k$ is the ordinary weighted average of $\log a_m$ and $\log a_n$, because $\log a_j$ is an arithmetic sequence.)
